\documentclass[aspectratio=169]{beamer}

\usetheme{Madrid}
\usecolortheme{default}
\setbeamertemplate{navigation symbols}{}

\usepackage[T1]{fontenc}
\usepackage{lmodern}
\usepackage{booktabs}
\usepackage{listings}

\lstset{
  basicstyle=\ttfamily\small,
  keywordstyle=\bfseries,
  columns=fullflexible,
  frame=single,
}

\title{Presentation Title}% TODO: title
\subtitle{Course or Conference}% TODO: subtitle
\author{Presenter Name}% TODO: presenter
\institute{Department, University}% TODO: institution
\date{\today}

\begin{document}

\begin{frame}
  \titlepage
\end{frame}

\begin{frame}{Outline}
  \tableofcontents
\end{frame}

\section{Background}

\begin{frame}{Problem}
  \begin{itemize}
    \item State the problem in one line.
    \item Say why it matters.
    \pause
    \item Show the key idea last.
  \end{itemize}
\end{frame}

\section{Method}

\begin{frame}{Gauss--Legendre quadrature}
  With two nodes,
  \begin{equation}
    \int_{-1}^{1} f(x)\,dx \approx f\!\left(-\frac{1}{\sqrt{3}}\right)
                                 + f\!\left(\frac{1}{\sqrt{3}}\right),
    \label{eq:gauss}
  \end{equation}
  which is exact for polynomials of degree at most~3.
\end{frame}

% Any frame with lstlisting, verbatim or \verb needs [fragile].
\begin{frame}[fragile]{Code}
  \begin{lstlisting}[language=Python]
def trapezoid(f, a, b, n):
    h = (b - a) / n
    s = (f(a) + f(b)) / 2 + sum(f(a + i*h) for i in range(1, n))
    return h * s
  \end{lstlisting}
\end{frame}

\begin{frame}{Results}
  \centering
  \begin{tabular}{@{}lr@{}}
    \toprule
    Method & Error \\
    \midrule
    Trapezoid ($n = 4$) & $2.6 \times 10^{-2}$ \\
    Gauss ($2$ nodes)   & $1.2 \times 10^{-3}$ \\
    \bottomrule
  \end{tabular}
\end{frame}

\begin{frame}{Summary}
  \begin{block}{Take-home message}
    One sentence the audience should remember.
  \end{block}
\end{frame}

\end{document}
